\documentclass[11pt]{article}

\usepackage{fullpage}
\usepackage{amsmath}
\usepackage{amssymb}
\usepackage{amsfonts}
\usepackage{amsthm}
\usepackage{mathtools}
\usepackage{enumitem}
\usepackage[colorlinks,linkcolor=magenta,citecolor=blue]{hyperref}

\newcommand{\Rbar}{\bar R}
\newcommand{\Om}{\Omega}
\newcommand{\E}{\mathbb E}
\newcommand{\Var}{\operatorname{Var}}
\newcommand{\tr}{\operatorname{tr}}
\newcommand{\diag}{\operatorname{diag}}
\newcommand{\Unif}{\operatorname{Unif}}
\newcommand{\IG}{\operatorname{IG}}
\newcommand{\Gam}{\operatorname{Gamma}}
\newcommand{\Normal}{\mathcal N}

\newtheorem{lemma}{Lemma}

\title{Better approximation to the susie-love model}
\author{}
\date{}

\begin{document}
\maketitle


\section{TO-DO}
Understand the relationship between 
\begin{enumerate}
\item PIP only depends on $x$ and the diagonal of $R$.
\item $p(x_{-j} | x_j, \beta, \gamma = j)$ does not depend on $\beta$.
\end{enumerate}
Look Yuxin's thesis and read the SI of susie-rss paper more carefully.

Why does having $\beta^2$ in the variance make sense? Of course it follows from the IW or Wi model, but intuitively why does it make sense?

Is there another model that arises from the conditional distribution:
\begin{equation}\label{eq:conditional-dist-Matthew}
p(x_{-j}\mid x_j, \beta, \gamma=j) = 
\Normal\left(
x_j \bar R_{j, -j} ,\,
\left(\frac{1}{N}+ \dfrac{x_j^2}{\nu_0 + 3}\right)S_j
\right)?
\end{equation}

\section{Some new models with better approximation to susie-love}
Recall our original SER model from what is called `susie-love':
\begin{align}\label{eq:likelihood-x-R-mean-var} 
\begin{cases}
x | R, \beta, \gamma \sim \mathcal{N}(\beta R_{\gamma}, R / N)\\
R \sim \mathcal{IW}(\nu_0 \bar R, \nu_0 + J + 1).
\end{cases}
\end{align}
Note that the latent matrix $R$ appears in both mean and variance of the likelihood in~\eqref{eq:likelihood-x-R-mean-var}. This model is very natural for unknown in-sample LD $R$, but is hard to make inferences. Last time we approximated $R$ in the variance to be $\bar R$, which makes the inference much easier but might eventually lead to the loss of the ``irrelevance of the null SNPs'' property. In principle, it makes sense that when $L = 1$, PIP of SNP j calculation should only depend on $x_j$ but not on other $x_k$ and $\bar R_{jk}$ (``irrelevance of the null SNPs'' property for SER).

I think this original model is actually feasible for inference. Consider the following equivalent way to represent it.

\subsection{Model A (equivalent to the original SER-love model)} Let's write down $p(x | R, \beta, \gamma = 1)$, other cases of $\gamma$ are similar. Let 
$$d_1 = R_{11} \in \mathbb{R}_{+}, \quad z_1 = \dfrac{R_{-1, 1}}{d_1} \in \mathbb{R}^{J-1},$$
and 
$$S_1 = R_{-1, -1} - d_1 z_1 z_1^\top $$
be the Schur's complement of $R_{11}$ in $R$. Hence, the latent $R$ is represented as
\begin{equation}\label{eq:R-block}
R = \begin{pmatrix}
d_1 & d_1 z_1^\top \\
d_1 z_1 & S_1 + d_1 z_1 z_1^\top
\end{pmatrix}.
\end{equation}
Similarly, denote
$$\bar S_1 = \bar R_{-1, -1} - \bar R_{-1, 1} \bar R_{-1, 1}^\top $$
be the Schur's complement of $\bar R_{11}$ in $\bar R$, and $\mu_1 =  \bar R_{-1, 1}$, we can write
\begin{equation}\label{eq:R-bar-block}
\bar R = \begin{pmatrix}
1 &\mu_1^\top \\
\mu_1 & \bar S_1 + \mu_1 \mu_1^\top
\end{pmatrix}.
\end{equation}

Properties of IW distribution (together with the fact that $\bar R_{11} = 1$)
imply, using the shape--scale parameterization for the inverse-gamma
distribution,
\begin{equation}
\begin{cases}
d_1 \sim \IG\left(\dfrac{\nu_0 + 2}{2}, \dfrac{\nu_0}{2}\right) \\
z_1 | S_1 \sim \mathcal{N}\left(\bar R_{-1, 1}, \dfrac{S_1}{\nu_0}\right)\\
S_1 \sim \mathcal{IW}(\nu_0 \bar S_1, \nu_0 + J + 1)
\end{cases}
\end{equation}
The original SER version of Susie-love (with $\gamma = 1$; similar for other $\gamma$) becomes:
\begin{align}\label{eq:SER-model-A} 
\begin{cases}
x | \beta, \gamma = 1, z_1, S_1 \sim \mathcal{N}\left(\beta 
\begin{pmatrix}
d_1  \\
d_1 z_1
\end{pmatrix}, 
\dfrac{1}{N} \begin{pmatrix}
d_1 & d_1 z_1^\top \\
d_1 z_1 & S_1 + d_1 z_1 z_1^\top
\end{pmatrix}\right)\\
d_1 \sim \IG\left(\dfrac{\nu_0 + 2}{2}, \dfrac{\nu_0}{2}\right) \\
z_1 | S_1 \sim \mathcal{N}\left(\mu_1, \dfrac{S_1}{\nu_0}\right)\\
S_1 \sim \mathcal{IW}(\nu_0 \bar S_1, \nu_0 + J + 1)\\
\beta \sim \mathcal{N}(0, \sigma_0^2)
\end{cases}
\end{align}

\paragraph{Deriving $p(x\mid \gamma=1)$.}
Write
\[
x=\begin{pmatrix}x_1\\x_{-1}\end{pmatrix}.
\]
\paragraph{Step 1: likelihood conditional on all latent variables.}
From the block normal distribution in~\eqref{eq:SER-model-A},
\[
p(x\mid \gamma=1,\beta,d_1,z_1,S_1)
=
\Normal_J\left(
x\mid
\beta
\begin{pmatrix}
d_1\\ d_1z_1
\end{pmatrix},
\frac{1}{N}
\begin{pmatrix}
d_1 & d_1z_1^\top\\
d_1z_1 & S_1+d_1z_1z_1^\top
\end{pmatrix}
\right).
\]
Equivalently, using the conditional distribution of a block normal,
\[
p(x\mid \gamma=1,\beta,d_1,z_1,S_1)
=
\Normal\left(x_1\mid \beta d_1,\frac{d_1}{N}\right)
\Normal_{J-1}\left(x_{-1}\mid x_1z_1,\frac{S_1}{N}\right).
\]
The key simplification is that, \textbf{after conditioning on $x_1$, the conditional
mean of $x_{-1}$ is $x_1z_1$ and no longer contains $\beta$ or $d_1$} (relevant to our earlier discussion).

\paragraph{Step 2: marginalize over $z_1\mid S_1$.}
Since
\[
z_1\mid S_1\sim \Normal_{J-1}\left(\mu_1,\frac{S_1}{\nu_0}\right)
\]
and
\[
x_{-1}\mid x_1,z_1,S_1
\sim
\Normal_{J-1}\left(x_1z_1,\frac{S_1}{N}\right),
\]
we get
\begin{equation}\label{eq:conditional-dist-susie-love}
x_{-1}\mid x_1,S_1,\gamma=1
\sim
\Normal_{J-1}\left(
x_1\mu_1,
\left(\frac{1}{N}+\frac{x_1^2}{\nu_0}\right)S_1
\right).
\end{equation}
Therefore
\[
p(x\mid \gamma=1,\beta,d_1,S_1)
=
\Normal\left(x_1\mid \beta d_1,\frac{d_1}{N}\right)
\Normal_{J-1}\left(
x_{-1}\mid x_1\mu_1,
\left(\frac{1}{N}+\frac{x_1^2}{\nu_0}\right)S_1
\right).
\]
What I found so interesting is that the conditional distribution $p(x_{-1}\mid x_1,S_1,\gamma=1)$ resembles what was suggested by Matthew last time (cf.~\eqref{eq:conditional-dist-Matthew}).

\paragraph{Step 3: marginalize over $S_1$.}
Because $S_1\sim \mathcal{IW}_{J-1}(\nu_0\bar S_1,\nu_0+J+1)$, integrating
$S_1$ gives a multivariate $t$ distribution with $\nu_0+3$ degrees of freedom:
\[
x_{-1}\mid x_1,\gamma=1
\sim
t_{J-1,\nu_0+3}\left(
x_1\mu_1,\,
\frac{\nu_0}{\nu_0+3}\left(\frac{1}{N}+\frac{x_1^2}{\nu_0}\right)\bar S_1
\right),
\]
where the second argument is the multivariate $t$ scale matrix. Thus
\[
p(x\mid \gamma=1,\beta,d_1)
=
\Normal\left(x_1\mid \beta d_1,\frac{d_1}{N}\right)
t_{J-1,\nu_0+3}\left(
x_{-1}\mid x_1\mu_1,\,
\frac{\nu_0}{\nu_0+3}\left(\frac{1}{N}+\frac{x_1^2}{\nu_0}\right)\bar S_1
\right).
\]
Equivalently, the $t$ factor is proportional to
\[
|\bar S_1|^{-1/2}
\left(\frac{1}{N}+\frac{x_1^2}{\nu_0}\right)^{-(J-1)/2}
\left[
1+
\frac{
(x_{-1}-x_1\mu_1)^\top(\nu_0\bar S_1)^{-1}
(x_{-1}-x_1\mu_1)
}{\frac{1}{N}+\frac{x_1^2}{\nu_0}}
\right]^{-(\nu_0+J+2)/2}.
\]
Using the block inverse formula (notice the block form of $\bar R$~\eqref{eq:R-bar-block}),
\[
x^\top\bar R^{-1}x
=
x_1^2+
(x_{-1}-x_1\mu_1)^\top\bar S_1^{-1}
(x_{-1}-x_1\mu_1),
\]
so
\[
1+
\frac{
(x_{-1}-x_1\mu_1)^\top(\nu_0\bar S_1)^{-1}
(x_{-1}-x_1\mu_1)
}{\frac{1}{N}+\frac{x_1^2}{\nu_0}}
=
\frac{x^\top\bar R^{-1}x+\nu_0/N}{x_1^2+\nu_0/N}.
\]
Therefore
\begin{equation}\label{eq:T_1}
T_1(x)
\propto
|\bar S_1|^{-1/2}
\left(x_1^2+\frac{\nu_0}{N}\right)^{(\nu_0+3)/2}
\left(x^\top\bar R^{-1}x+\frac{\nu_0}{N}\right)^{-(\nu_0+J+2)/2}.
\end{equation}
When comparing different values of $\gamma$, the final factor is common across
SNPs. Also, since $\bar R_{11}=1$, $|\bar R|=|\bar S_1|$; analogously,
$|\bar S_j|=|\bar R|$ for any SNP $j$ with $\bar R_{jj}=1$.



\paragraph{Step 4: marginalize over $\beta$ and $d_1$.}
The preceding display separates the problem into an $x_{-1}\mid x_1$ factor
and an $x_1$ factor. Define
\[
T_1(x)=
t_{J-1,\nu_0+3}\left(
x_{-1}\mid x_1\mu_1,\,
\frac{\nu_0}{\nu_0+3}\left(\frac{1}{N}+\frac{x_1^2}{\nu_0}\right)\bar S_1
\right).
\]
If we integrate over $\beta$, then
\[
x_1\mid d_1,\gamma=1
\sim
\Normal\left(0,\frac{d_1}{N}+d_1^2\sigma_0^2\right),
\]
and hence (finally integrating over $d_1$):

\[
\boxed{
p(x\mid \gamma=1)
=
T_1(x)
\int_0^\infty
\Normal\left(x_1\mid 0,\frac{d_1}{N}+d_1^2\sigma_0^2\right)
\IG\left(d_1;\frac{\nu_0+2}{2},\frac{\nu_0}{2}\right)\,d d_1
}.
\]
This is a one-dimensional integral over $d_1>0$. I do not see a standard
closed form for it because the variance is $d_1/N+d_1^2\sigma_0^2$. Replacing $T_1$ with its expression in~\eqref{eq:T_1}, 

\[
\begin{aligned}
p(x\mid \gamma=1)
\propto\;& |\bar R|^{-1/2}
\left(x_1^2+\frac{\nu_0}{N}\right)^{(\nu_0+3)/2}
\left(x^\top\bar R^{-1}x+\frac{\nu_0}{N}\right)^{-(\nu_0+J+2)/2} \\
&\quad\times
\int_0^\infty
\Normal\left(x_1\mid 0,\frac{d_1}{N}+d_1^2\sigma_0^2\right)\,d \IG(d_1)
\end{aligned}
\]
Since $|\bar R|$ and $x^\top\bar R^{-1}x$ do not depend on which SNP is chosen,
both factors cancel in the PIP normalization. Also,
\[
\left(x_j^2+\frac{\nu_0}{N}\right)^{(\nu_0+3)/2}
=
\left(\frac{\nu_0}{N}\right)^{(\nu_0+3)/2}
\left(1+\frac{Nx_j^2}{\nu_0}\right)^{(\nu_0+3)/2},
\]
and the first factor is common across SNPs. Hence,
\[
\boxed{
p(\gamma=j\mid x)
\propto
\left(1+\frac{Nx_j^2}{\nu_0}\right)^{(\nu_0+3)/2}
\int_0^\infty
\Normal\left(x_j\mid 0,\frac{d_j}{N}+d_j^2\sigma_0^2\right)\,d \IG(d_j)
}.
\]
\textbf{Conclusion.} In the original SER-love model, the marginal likelihood
depends on $x^\top\bar R^{-1}x$, but this dependence is common across SNPs and
cancels when computing the PIP. The resulting PIP weight depends on SNP $j$
only through $x_j$.

\subsection{Model B} 
Simplify $d_1 \equiv 1$ to avoid intractable integral (and we know the diagonal of the in-sample LD matrix is 1 anyway), we have the following model (model B):
\begin{align}\label{eq:SER-model-B} 
\begin{cases}
x | \beta, \gamma = 1, z_1, S_1 \sim \mathcal{N}\left(\beta 
\begin{pmatrix}
1  \\
z_1
\end{pmatrix}, 
\dfrac{1}{N} \begin{pmatrix}
1 & z_1^\top \\
z_1 & S_1 + z_1 z_1^\top
\end{pmatrix}\right)\\
z_1 | S_1 \sim \mathcal{N}\left(\mu_1, \dfrac{S_1}{\nu_0}\right)\\
S_1 \sim \mathcal{IW}(\nu_0 \bar S_1, \nu_0 + J + 1)\\
\beta \sim \mathcal{N}(0, \sigma_0^2)
\end{cases}
\end{align}

This model leads to
\[
\boxed{
p(x\mid\gamma=1)
\propto
|\bar R|^{-1/2}
\left(1+\frac{Nx_j^2}{\nu_0}\right)^{(\nu_0+3)/2}
\left(x^\top\bar R^{-1}x+\frac{\nu_0}{N}\right)^{-(\nu_0+J+2)/2}
\Normal\left(x_1\mid 0,\frac{1}{N} + \sigma_0^2\right)
}.
\]
Since $|\bar R|$ and $x^\top\bar R^{-1}x$ do not depend on which SNP is chosen,
the PIP weights simplify further. Dropping the common factors gives
\begin{equation}\label{eq:PIP-model-B}
\boxed{
p(\gamma=j \mid x)
\propto
\left(1+\frac{Nx_j^2}{\nu_0}\right)^{(\nu_0+3)/2}
\Normal\left(x_j \mid 0,\frac{1}{N} + \sigma_0^2\right)
}.
\end{equation}

\paragraph{Compare to SER-RSS's PIP.} Recall SER-RSS's PIP:
\begin{equation}\label{eq:PIP-susie-RSS}
\tilde{p}(\gamma=j\mid x)
\propto \exp\left(\dfrac{N\sigma_0^{2}}{N\sigma_0^{2} + 1} \dfrac{N x_j^2}{2} \right)
\end{equation}
We now show that~\eqref{eq:PIP-model-B} tends to~\eqref{eq:PIP-susie-RSS}
as $\nu_0\to\infty$. Since
\[
\left(1+\frac{Nx_j^2}{\nu_0}\right)^{(\nu_0+3)/2} = \left(1+ \dfrac{2}{\nu_0 + 3} \times \dfrac{\nu_0 + 3}{\nu_0} \frac{Nx_j^2}{2}\right)^{(\nu_0+3)/2} 
\asymp \exp\left\{\frac{\nu_0 + 3}{\nu_0} \dfrac{Nx_j^2}{2}\right\}
\]
Also,
\[
\Normal\left(x_j\mid 0,\frac{1}{N}+\sigma_0^2\right)
\propto
\exp\left\{
-\frac{x_j^2}{2(1/N+\sigma_0^2)}
\right\}
=
\exp\left\{
-\frac{1}{1+N\sigma_0^2} \frac{N x_j^2}{2}
\right\}.
\]
Therefore the limiting weight in~\eqref{eq:PIP-model-B} is asymptotically (as $\nu_0$ gets large)
\[
\exp\left\{\left(\frac{\nu_0 + 3}{\nu_0} - \dfrac{1}{N\sigma_0^2 + 1}\right) \dfrac{Nx_j^2}{2}\right\} = \exp\left\{\left(\dfrac{N\sigma_0^2}{N\sigma_0^2 + 1} - \dfrac{1}{\nu_0 + 3}\right) \dfrac{Nx_j^2}{2}\right\} \approx \exp\left\{\dfrac{N\sigma_0^2}{N\sigma_0^2 + 1} \dfrac{Nx_j^2}{2}\right\}
\]
which is exactly the SER-RSS PIP weight in~\eqref{eq:PIP-susie-RSS}. Hence, for finite $\nu_0$, PIP under SER-love model (B) is slightly more diffused than susie-rss. 

\paragraph{Inference under model B~\eqref{eq:SER-model-B}.}
For a general SNP $j$, write
\[
\mu_j=\bar R_{-j,j},
\qquad
\bar S_j=\bar R_{-j,-j}-\mu_j\mu_j^\top,
\qquad
\kappa_j=\nu_0+Nx_j^2.
\]
Under model B, conditional on $\gamma=j$, the likelihood factorizes as
\[
p(x\mid \beta,z_j,S_j,\gamma=j)
=
\Normal\left(x_j\mid \beta,\frac{1}{N}\right)
\Normal_{J-1}\left(x_{-j}\mid x_jz_j,\frac{S_j}{N}\right).
\]
Thus $\beta$ is conditionally independent of $(z_j,S_j)$ given $x$ and
$\gamma=j$, and
\[
p(\beta,z_j,S_j\mid x,\gamma=j)
=
p(\beta\mid x_j,\gamma=j)\,
p(z_j,S_j\mid x,\gamma=j).
\]

$\bullet$ The posterior for $\beta$ is the usual normal-normal update:
\[
\boxed{
\beta\mid x,\gamma=j
\sim
\Normal\left(
\frac{N}{N+\sigma_0^{-2}}x_j,\,
\frac{1}{N+\sigma_0^{-2}}
\right).
}
\]

$\bullet$ For the LD-column variables, first condition on $S_j$. Combining
\[
z_j\mid S_j\sim \Normal_{J-1}\left(\mu_j,\frac{S_j}{\nu_0}\right),
\qquad
x_{-j}\mid x_j,z_j,S_j,\gamma=j
\sim
\Normal_{J-1}\left(x_jz_j,\frac{S_j}{N}\right)
\]
gives
\[
\boxed{
z_j\mid S_j,x,\gamma=j
\sim
\Normal_{J-1}\left(
\frac{\nu_0\mu_j+Nx_jx_{-j}}{\nu_0+Nx_j^2},\,
\frac{S_j}{\nu_0+Nx_j^2}
\right).
}
\]

$\bullet$ Integrating out $z_j$ gives
\[
x_{-j}\mid x_j,S_j,\gamma=j
\sim
\Normal_{J-1}\left(
x_j\mu_j,\,
\left(\frac{1}{N}+\frac{x_j^2}{\nu_0}\right)S_j
\right).
\]
Therefore the posterior for $S_j$ is inverse-Wishart:
\[
\boxed{
S_j\mid x,\gamma=j
\sim
\mathcal{IW}\left(
\nu_0\bar S_j+
\frac{N\nu_0}{\nu_0+Nx_j^2}
(x_{-j}-x_j\mu_j)(x_{-j}-x_j\mu_j)^\top,\,
\nu_0+J+2
\right).
}
\]
Equivalently, the joint posterior of $(z_j,S_j)$ can be represented by first
drawing $S_j$ from this inverse-Wishart distribution and then drawing $z_j$
from the conditional normal distribution above. Finally, the full posterior
over $\gamma$ uses the PIP weights in~\eqref{eq:PIP-model-B}, normalized over
$j=1,\ldots,J$.

\paragraph{Closed-form ELBO.}
Because the conditional posteriors above are exact, the optimized ELBO for
model B is the exact marginal log likelihood. Let
\[
V_\beta=\frac{1}{N}+\sigma_0^2,
\qquad
q=x^\top\bar R^{-1}x.
\]
For a uniform prior on $\gamma$, the optimized ELBO is
\[
\mathcal L(\nu_0,\sigma_0^2)
=
\log\left\{
\frac{1}{J}\sum_{j=1}^J \exp(\ell_j)
\right\},
\]
where the component log marginal likelihood can be written as
\begin{align*}
\ell_j
=\;&
-\frac{1}{2}\log(2\pi V_\beta)
-\frac{x_j^2}{2V_\beta}
+\log\Gamma\left(\frac{\nu_0+J+2}{2}\right)
-\log\Gamma\left(\frac{\nu_0+3}{2}\right) \\
&\quad
-\frac{J-1}{2}\log\pi
-\frac{1}{2}\log|\bar S_j|
+\frac{\nu_0+3}{2}\log\left(x_j^2+\frac{\nu_0}{N}\right)
-\frac{\nu_0+J+2}{2}\log\left(q+\frac{\nu_0}{N}\right).
\end{align*}
This is equivalent to the marginal likelihood derived above, but unlike the
PIP formula it keeps the terms that are common across SNPs. Those common terms
do not affect $\pi_j=p(\gamma=j\mid x)$ for fixed hyperparameters, but they do
matter when learning $\nu_0$.

\paragraph{Empirical Bayes learning of $\nu_0$ and $\sigma_0^2$.}
The empirical Bayes estimates can be obtained by maximizing the closed-form
objective
\[
(\hat\nu_0,\hat\sigma_0^2)
=
\arg\max_{\nu_0>0,\;\sigma_0^2\ge 0}
\mathcal L(\nu_0,\sigma_0^2).
\]
After choosing $(\hat\nu_0,\hat\sigma_0^2)$, the posterior inclusion
probabilities are
\[
\pi_j
=
\frac{\exp(\ell_j)}
{\sum_{k=1}^J \exp(\ell_k)}
\bigg|_{\nu_0=\hat\nu_0,\;\sigma_0^2=\hat\sigma_0^2}.
\]
Equivalently, for computing PIPs at fixed hyperparameters one may use the
simplified weights in~\eqref{eq:PIP-model-B}; for empirical Bayes optimization
one should use $\mathcal L(\nu_0,\sigma_0^2)$ above, because it retains the
$\nu_0$-dependent normalizing constants and the common
$\log(q+\nu_0/N)$ term.

\paragraph{EM view of empirical Bayes.}
Another way to use the fact that the posterior is exact is to write the usual
ELBO identity. Let
\[
\theta=(\gamma,\beta,z_\gamma,S_\gamma),
\qquad
q(\theta)=p(\theta\mid x;\nu_0^{\mathrm{old}},\sigma_{0,\mathrm{old}}^2).
\]
Strictly speaking,
\[
\log p(x;\nu_0,\sigma_0^2)
\ge
\E_q\{\log p(x\mid\theta)\}
-\operatorname{KL}\{q(\theta)\,\|\,p(\theta;\nu_0,\sigma_0^2)\},
\]
with equality when $q$ is the posterior under the same hyperparameters. In an
EM algorithm, the E-step fixes $q$ at the current hyperparameters and the
M-step maximizes
\[
Q(\nu_0,\sigma_0^2)
=
\E_q\{\log p(\theta;\nu_0,\sigma_0^2)\},
\]
because the likelihood term does not involve $\nu_0$ or $\sigma_0^2$ in model B.

This gives a closed-form update for $\sigma_0^2$. If
\[
\pi_j=q(\gamma=j),
\qquad
m_{\beta j}=\frac{N}{N+\sigma_{0,\mathrm{old}}^{-2}}x_j,
\qquad
v_\beta=\frac{1}{N+\sigma_{0,\mathrm{old}}^{-2}},
\]
then
\[
\boxed{
\sigma_{0,\mathrm{new}}^2
=
\sum_{j=1}^J \pi_j\left(m_{\beta j}^2+v_\beta\right).
}
\]

For $\nu_0$, the same EM view reduces the update to a one-dimensional
optimization, but not to a closed-form expression. Let $p=J-1$ and write the
current posterior of $S_j$ as
\[
S_j\mid x,\gamma=j
\sim
\mathcal{IW}_p(\Psi_j,m_j),
\qquad
m_j=\nu_0^{\mathrm{old}}+J+2,
\]
where
\[
\Psi_j
=
\nu_0^{\mathrm{old}}\bar S_j+
\frac{N\nu_0^{\mathrm{old}}}{\nu_0^{\mathrm{old}}+Nx_j^2}
(x_{-j}-x_j\mu_j)(x_{-j}-x_j\mu_j)^\top.
\]
Also let
\[
\bar z_j
=
\frac{\nu_0^{\mathrm{old}}\mu_j+Nx_jx_{-j}}
{\nu_0^{\mathrm{old}}+Nx_j^2},
\qquad
\kappa_j^{\mathrm{old}}=\nu_0^{\mathrm{old}}+Nx_j^2.
\]
These quantities can be evaluated without forming or inverting the
$(J-1)\times(J-1)$ matrices $\Psi_j$. Let
\[
\nu_\star=\nu_0^{\mathrm{old}},
\qquad
q_R=x^\top\bar R^{-1}x,
\qquad
\delta_j=x_{-j}-x_j\mu_j.
\]
Then
\[
\Psi_j
=
\nu_\star\bar S_j+
\frac{N\nu_\star}{\kappa_j^{\mathrm{old}}}\delta_j\delta_j^\top.
\]
By the Woodbury identity,
\[
\Psi_j^{-1}
=
\frac{1}{\nu_\star}\bar S_j^{-1}
-
\frac{N}{\nu_\star(\nu_\star+Nq_R)}
\bar S_j^{-1}\delta_j\delta_j^\top\bar S_j^{-1},
\]
because
\[
1+\frac{N}{\kappa_j^{\mathrm{old}}}(q_R-x_j^2)
=
\frac{\nu_\star+Nq_R}{\kappa_j^{\mathrm{old}}}.
\]
The determinant lemma gives
\[
\log|\Psi_j|
=
(J-1)\log\nu_\star+\log|\bar S_j|
+\log\left(\frac{\nu_\star+Nq_R}{\kappa_j^{\mathrm{old}}}\right).
\]
Since $\bar R_{jj}=1$, $\log|\bar S_j|=\log|\bar R|$, so this is also common
to compute once.

The needed posterior expectations are
\[
A_j=\E_q\{\tr(\bar S_jS_j^{-1})\}
=
m_j
\left\{
\frac{J-1}{\nu_\star}
-
\frac{N(q_R-x_j^2)}{\nu_\star(\nu_\star+Nq_R)}
\right\},
\]
\[
B_j=\E_q\{(z_j-\mu_j)^\top S_j^{-1}(z_j-\mu_j)\}
=
m_j
\frac{N^2x_j^2(q_R-x_j^2)}
{\nu_\star\kappa_j^{\mathrm{old}}(\nu_\star+Nq_R)}
+\frac{p}{\kappa_j^{\mathrm{old}}},
\]
and
\[
C_j=\E_q\{\log|S_j|\}
=
\;(J-1)\log\nu_\star+\log|\bar R|
+\log\left(\frac{\nu_\star+Nq_R}{\kappa_j^{\mathrm{old}}}\right)
-\sum_{r=1}^p \psi\left(\frac{m_j+1-r}{2}\right)
-p\log 2.
\]
Thus all SNP-specific parts of $A_j,B_j,C_j$ depend on $j$ only through
$x_j^2$ and $q_R-x_j^2$. In implementation, one can compute $q_R$ and
$\log|\bar R|$ once, then evaluate the remaining terms with scalar operations
for each SNP.
Up to constants that do not depend on $\nu_0$, the EM objective for $\nu_0$ is
\begin{align*}
Q_\nu(\nu_0)
=
\sum_{j=1}^J \pi_j
\Bigg[
&\frac{p}{2}\log\nu_0
+\frac{\nu_0+J+1}{2}\log|\nu_0\bar S_j|
-\frac{(\nu_0+J+1)p}{2}\log 2 \\
&-\log\Gamma_p\left(\frac{\nu_0+J+1}{2}\right)
-\frac{\nu_0+2J+2}{2}C_j
-\frac{\nu_0}{2}(A_j+B_j)
\Bigg].
\end{align*}
Thus
\[
\boxed{
\nu_{0,\mathrm{new}}=\arg\max_{\nu_0>0} Q_\nu(\nu_0).
}
\]
So the exact-posterior ELBO gives a particularly nice update for
$\sigma_0^2$, and a stable scalar optimization for $\nu_0$. For $\nu_0$, the
closed-form marginal likelihood $\mathcal L(\nu_0,\sigma_0^2)$ above is often
even simpler to optimize directly, but the EM form is useful for deriving a
monotone coordinate-ascent algorithm.

\subsection{Other representations of Model B}
Recall Model B:

\begin{align}\label{eq:SER-model-B-recall} 
\begin{cases}
x | \beta, \gamma = 1, z_1, S_1 \sim \mathcal{N}\left(\beta 
\begin{pmatrix}
1  \\
z_1
\end{pmatrix}, 
\dfrac{1}{N} \begin{pmatrix}
1 & z_1^\top \\
z_1 & S_1 + z_1 z_1^\top
\end{pmatrix}\right)\\
z_1 | S_1 \sim \mathcal{N}\left(\mu_1, \dfrac{S_1}{\nu_0}\right)\\
S_1 \sim \mathcal{IW}(\nu_0 \bar S_1, \nu_0 + J + 1)\\
\beta \sim \mathcal{N}(0, \sigma_0^2)
\end{cases}
\end{align}

or equivalently,

\begin{align}\label{eq:model-Bp}
\begin{cases}
x_1\mid \beta,\gamma=1
\sim
\Normal\left(\beta,\dfrac{1}{N}\right),\\
x_{-1}\mid x_1,z_1,\gamma=1
\sim
\Normal_{J-1}\left(
x_1 z_1,
\dfrac{1}{N}S_1
\right),\\
z_1 | S_1 \sim \mathcal{N}\left(\mu_1, \dfrac{S_1}{\nu_0}\right),\\
S_1\sim\mathcal{IW}(\nu_0\bar S_1,\nu_0+J+1),\\
\beta\sim\Normal(0,\sigma_0^2).
\end{cases}
\end{align}

Here are three other equivalent representations of model B for $\gamma=1$. Write
\[
\mu_1=\bar R_{-1,1},
\qquad
\bar S_1=\bar R_{-1,-1}-\mu_1\mu_1^\top.
\]
\paragraph{1. After integrating out $z_1\mid S_1$.}
\begin{align}\label{eq:model-B1}
\begin{cases}
x_1\mid \beta,\gamma=1
\sim
\Normal\left(\beta,\dfrac{1}{N}\right),\\
x_{-1}\mid x_1,S_1,\gamma=1
\sim
\Normal_{J-1}\left(
x_1\mu_1,
\left(\dfrac{1}{N}+\dfrac{x_1^2}{\nu_0}\right)S_1
\right),\\
S_1\sim\mathcal{IW}(\nu_0\bar S_1,\nu_0+J+1),\\
\beta\sim\Normal(0,\sigma_0^2).
\end{cases}
\end{align}

\paragraph{2. After further integrating out $S_1$.}
\begin{align}\label{eq:model-B2}
\begin{cases}
x_1\mid \beta,\gamma=1
\sim
\Normal\left(\beta,\dfrac{1}{N}\right),\\
x_{-1}\mid x_1,\gamma=1
\sim
t_{J-1,\nu_0+3}\left(
x_1\mu_1,
\dfrac{\nu_0}{\nu_0+3}
\left(\dfrac{1}{N}+\dfrac{x_1^2}{\nu_0}\right)\bar S_1
\right),\\
\beta\sim\Normal(0,\sigma_0^2).
\end{cases}
\end{align}

\paragraph{3. Using the normal--Gamma mixture for the $t$ distribution.}
With shape--rate parameterization,
Conditional on $\omega_1$:
\begin{align}\label{eq:model-B3p}
\begin{cases}
x_1\mid \beta,\gamma=1
\sim
\Normal\left(\beta,\dfrac{1}{N}\right),\\
x_{-1}\mid x_1,\omega_1,\gamma=1
\sim
\Normal_{J-1}\left(
x_1\mu_1,
\dfrac{1}{\omega_1}
\left(\dfrac{1}{N}+\dfrac{x_1^2}{\nu_0}\right)\bar S_1
\right),\\
\omega_1\sim\Gam\left(\dfrac{\nu_0+3}{2},\dfrac{\nu_0}{2}\right),\\
\beta\sim\Normal(0,\sigma_0^2).
\end{cases}
\end{align}

\subsection{Try to achieve common noise}

I am thinking about another representation of model B~\eqref{eq:SER-model-B}
that will be useful in extending it to a Susie-like model. In particular, it
would be useful to have a model with a common noise among all choices of $\gamma = j$.

\paragraph{An exact two-error view of model B.}
For $\gamma=1$, model B can be viewed as a SER model with two additional
errors, but the variance error is not an independent additive Gaussian error.
Let $A_1=A_1(S_1)$ be any matrix such that
\[
A_1\bar S_1A_1^\top=S_1.
\]
Let
\[
e_1\sim\Normal\left(0,\frac{1}{N}\right),
\qquad
e_{-1}\sim\Normal_{J-1}\left(0,\frac{\bar S_1}{N}\right),
\qquad
h_1\sim\Normal_{J-1}\left(0,\frac{\bar S_1}{\nu_0}\right),
\]
mutually independent conditional on $S_1$, and set
\begin{align}\label{eq:model-B-two-error-view}
\begin{cases}
x_1=\beta+e_1,\\
x_{-1}=x_1\mu_1+A_1e_{-1}+x_1A_1h_1,\\
S_1\sim\mathcal{IW}(\nu_0\bar S_1,\nu_0+J+1),\\
\beta\sim\Normal(0,\sigma_0^2).
\end{cases}
\end{align}
The term $x_1A_1h_1$ is the error in the mean/LD column, because
$A_1h_1\mid S_1\sim\Normal_{J-1}(0,S_1/\nu_0)$. The term $A_1e_{-1}$ is the
variance error: it transforms the SER conditional noise
$e_{-1}\sim\Normal_{J-1}(0,\bar S_1/N)$ into
\[
A_1e_{-1}\mid S_1\sim\Normal_{J-1}\left(0,\frac{S_1}{N}\right).
\]
Consequently, after integrating out $e_{-1}$ and $h_1$,
\[
x_{-1}\mid x_1,S_1,\gamma=1
\sim
\Normal_{J-1}\left(
x_1\mu_1,
\left(\frac{1}{N}+\frac{x_1^2}{\nu_0}\right)S_1
\right),
\]
which is exactly \eqref{eq:model-B1} in block form.

Equivalently, if
$x^{\mathrm{SER}}\mid\beta,\gamma=1\sim
\Normal_J(\beta\bar R_{\cdot 1},\bar R/N)$, then the conditional SER residual
in the non-causal coordinates is $e_{-1}=x_{-1}^{\mathrm{SER}}-x_1\mu_1$, and
\[
x
=
x^{\mathrm{SER}}
+u_1\left\{(A_1-I)e_{-1}+x_1A_1h_1\right\}.
\]
This is the desired ``SER plus two errors'' picture. The caveat is that
$(A_1-I)e_{-1}$ is coupled to the SER noise; it is not an independent
additive Gaussian error. This coupling is what avoids the impossible
requirement that $S_1-\bar S_1$ be positive semidefinite.

What is the distribution of $A_1$? It is not unique unless we choose a
particular square-root convention. If
\[
T_1=\bar S_1^{-1/2}S_1\bar S_1^{-1/2},
\]
then by affine equivariance of the inverse-Wishart distribution,
\[
T_1\sim
\mathcal{IW}_{J-1}\left(\nu_0 I,\nu_0+J+1\right).
\]
Every valid $A_1$ can be written as
\[
A_1=\bar S_1^{1/2}B_1\bar S_1^{-1/2},
\qquad
B_1B_1^\top=T_1.
\]
Thus one canonical choice is
\[
B_1=T_1^{1/2},
\qquad
A_1=\bar S_1^{1/2}T_1^{1/2}\bar S_1^{-1/2}.
\]
Under this convention, $A_1$ has the distribution induced by the symmetric
matrix square root of an isotropic inverse-Wishart random matrix. This does
not have a simpler standard named form, but it is concentrated around $I$ as
$\nu_0$ grows. Other choices, such as $B_1=T_1^{1/2}O_1$ for an orthogonal
matrix $O_1$, give the same $S_1$ but a different distribution for $A_1$.

\paragraph{Approximating the variance of $x_1$.}
Starting from \eqref{eq:model-B3p}, try the approximation
\[
x_1\mid \beta,\omega_1,\gamma=1
\sim
\Normal\left(\beta,\frac{1}{N\omega_1}\right),
\]
while keeping the conditional model for $x_{-1}$ given $x_1$ and $\omega_1$:
\begin{align}\label{eq:model-B3p-x1-scaled}
\begin{cases}
x_1\mid \beta,\omega_1,\gamma=1
\sim
\Normal\left(\beta,\dfrac{1}{N\omega_1}\right),\\
x_{-1}\mid x_1,\omega_1,\gamma=1
\sim
\Normal_{J-1}\left(
x_1\mu_1,
\dfrac{1}{\omega_1}
\left(\dfrac{1}{N}+\dfrac{x_1^2}{\nu_0}\right)\bar S_1
\right),\\
\omega_1\sim\Gam\left(\dfrac{\nu_0+3}{2},\dfrac{\nu_0}{2}\right),\\
\beta\sim\Normal(0,\sigma_0^2).
\end{cases}
\end{align}
This approximation has a clean whole-vector representation. Let
\[
\zeta_1\mid x_1,\omega_1,\gamma=1
\sim
\Normal_{J-1}\left(
0,
\dfrac{x_1^2}{\nu_0\omega_1}\bar S_1
\right).
\]
Then \eqref{eq:model-B3p-x1-scaled} is equivalent to
\begin{align}\label{eq:model-B3p-x1-scaled-vector}
\begin{cases}
x\mid \beta,\zeta_1,\omega_1,\gamma=1
\sim
\Normal_J\left(
\beta\bar R_{\cdot 1}+u_1(\zeta_1),
\dfrac{\bar R}{N\omega_1}
\right),\\
\zeta_1\mid x_1,\omega_1,\gamma=1
\sim
\Normal_{J-1}\left(
0,
\dfrac{x_1^2}{\nu_0\omega_1}\bar S_1
\right),\\
\omega_1\sim\Gam\left(\dfrac{\nu_0+3}{2},\dfrac{\nu_0}{2}\right),\\
\beta\sim\Normal(0,\sigma_0^2).
\end{cases}
\end{align}
Indeed, the conditional distribution induced by the common noise
$\bar R/(N\omega_1)$ is
\[
x_1\mid \beta,\omega_1,\gamma=1
\sim
\Normal\left(\beta,\frac{1}{N\omega_1}\right),
\qquad
x_{-1}\mid x_1,\zeta_1,\omega_1,\gamma=1
\sim
\Normal_{J-1}\left(
x_1\mu_1+\zeta_1,
\frac{\bar S_1}{N\omega_1}
\right).
\]
Integrating out $\zeta_1$ gives the second line of
\eqref{eq:model-B3p-x1-scaled}.

Thus this approximation achieves a common-noise representation with common
noise $\bar R/(N\omega_1)$. This is not the same as fixed noise $\bar R/N$. However, if we are willing
to use the random common-noise scale $\bar R/(N\omega_1)$, then
\eqref{eq:model-B3p-x1-scaled-vector} is a valid representation.

\paragraph{PIP under the random common-noise model.}
For a general candidate SNP $j$, define
\[
\mu_j=\bar R_{-j,j},
\qquad
\bar S_j=\bar R_{-j,-j}-\mu_j\mu_j^\top,
\qquad
c_j=\frac{1}{N}+\frac{x_j^2}{\nu_0},
\qquad
q_R=x^\top\bar R^{-1}x.
\]
The block inverse formula gives
\[
(x_{-j}-x_j\mu_j)^\top\bar S_j^{-1}(x_{-j}-x_j\mu_j)
=q_R-x_j^2.
\]
Under \eqref{eq:model-B3p-x1-scaled}, after integrating out $\beta$ but
conditioning on $\omega_j$,
\[
x_j\mid \omega_j,\gamma=j
\sim
\Normal\left(0,\sigma_0^2+\frac{1}{N\omega_j}\right),
\]
and
\[
x_{-j}\mid x_j,\omega_j,\gamma=j
\sim
\Normal_{J-1}\left(
x_j\mu_j,
\frac{c_j}{\omega_j}\bar S_j
\right).
\]
Therefore
\begin{align}
p(x\mid \gamma=j)
=\;&
\int_0^\infty
\Normal\left(x_j;0,\sigma_0^2+\frac{1}{N\omega}\right)
\Normal_{J-1}\left(
x_{-j};x_j\mu_j,\frac{c_j}{\omega}\bar S_j
\right)
\notag\\
&\qquad\qquad\times
\frac{(\nu_0/2)^{(\nu_0+3)/2}}
{\Gamma((\nu_0+3)/2)}
\omega^{(\nu_0+3)/2-1}
\exp\left(-\frac{\nu_0\omega}{2}\right)
\;d\omega \notag\\
=\;&
(2\pi)^{-J/2}|\bar S_j|^{-1/2}
c_j^{-(J-1)/2}
\frac{(\nu_0/2)^{(\nu_0+3)/2}}
{\Gamma((\nu_0+3)/2)}
\notag\\
&\times
\int_0^\infty
\left(\sigma_0^2+\frac{1}{N\omega}\right)^{-1/2}
\omega^{(\nu_0+3)/2+(J-1)/2-1}
\notag\\
&\times
\exp\left\{
-\frac{x_j^2}{2(\sigma_0^2+1/(N\omega))}
-\frac{\omega(q_R-x_j^2)}{2c_j}
-\frac{\nu_0\omega}{2}
\right\}
\;d\omega .
\label{eq:common-noise-scaled-marginal}
\end{align}
Since $\bar R_{jj}=1$, $|\bar S_j|=|\bar R|$, so the determinant factor is
common across SNPs. Dropping constants common in $j$, the PIP weights are
\begin{align}\label{eq:PIP-common-noise-scaled}
w_j^{\mathrm{scaled}}
=\;&
c_j^{-(J-1)/2}
\int_0^\infty
\left(\sigma_0^2+\frac{1}{N\omega}\right)^{-1/2}
\omega^{(\nu_0+J)/2}
\notag\\
&\times
\exp\left\{
-\frac{x_j^2}{2(\sigma_0^2+1/(N\omega))}
-\frac{\omega(q_R-x_j^2)}{2c_j}
-\frac{\nu_0\omega}{2}
\right\}
\;d\omega,
\qquad
\pi_j=\frac{w_j^{\mathrm{scaled}}}{\sum_{k=1}^Jw_k^{\mathrm{scaled}}}.
\end{align}
Thus the random common-noise model is computationally a one-dimensional
integral per SNP, but it does not retain the exact irrelevance property. The
common statistic $q_R=x^\top\bar R^{-1}x$ appears inside the integral with the
SNP-specific coefficient $1/c_j$, so it cannot be factored out and cancelled
in the PIP normalization.

\paragraph{A variant that scales the mean of $x_1$.}
Another possibility, closer in spirit to model A, is to replace the first line
of \eqref{eq:model-B3p} by
\[
x_1\mid \beta,\omega_1,\gamma=1
\sim
\Normal\left(\frac{\beta}{\sqrt{\omega_1}},\frac{1}{N\omega_1}\right).
\]
Keeping the same conditional distribution of $x_{-1}$, the block model becomes
\begin{align}\label{eq:model-B3p-x1-mean-scaled}
\begin{cases}
x_1\mid \beta,\omega_1,\gamma=1
\sim
\Normal\left(\dfrac{\beta}{\sqrt{\omega_1}},\dfrac{1}{N\omega_1}\right),\\
x_{-1}\mid x_1,\omega_1,\gamma=1
\sim
\Normal_{J-1}\left(
x_1\mu_1,
\dfrac{1}{\omega_1}
\left(\dfrac{1}{N}+\dfrac{x_1^2}{\nu_0}\right)\bar S_1
\right),\\
\omega_1\sim\Gam\left(\dfrac{\nu_0+3}{2},\dfrac{\nu_0}{2}\right),\\
\beta\sim\Normal(0,\sigma_0^2).
\end{cases}
\end{align}
This has the common-noise representation
\begin{align}\label{eq:model-B3p-x1-mean-scaled-vector}
\begin{cases}
x\mid \beta,\zeta_1,\omega_1,\gamma=1
\sim
\Normal_J\left(
\dfrac{\beta}{\sqrt{\omega_1}}\bar R_{\cdot 1}+u_1(\zeta_1),
\dfrac{\bar R}{N\omega_1}
\right),\\
\zeta_1\mid x_1,\omega_1,\gamma=1
\sim
\Normal_{J-1}\left(
0,
\dfrac{x_1^2}{\nu_0\omega_1}\bar S_1
\right),\\
\omega_1\sim\Gam\left(\dfrac{\nu_0+3}{2},\dfrac{\nu_0}{2}\right),\\
\beta\sim\Normal(0,\sigma_0^2).
\end{cases}
\end{align}
The only difference from \eqref{eq:model-B3p-x1-scaled-vector} is that the
effect mean is also divided by $\sqrt{\omega_1}$. If we define
$\alpha_1=\beta/\sqrt{\omega_1}$, then
\[
\alpha_1\mid \omega_1\sim\Normal\left(0,\frac{\sigma_0^2}{\omega_1}\right),
\]
so the prior variance of the effective effect size has the same $1/\omega_1$
scaling as the common noise. This is the useful feature of the
$1/\sqrt{\omega_1}$ scaling.

The PIP calculation is almost the same as above. Conditional on $\omega_j$ and
after integrating out $\beta$,
\[
x_j\mid \omega_j,\gamma=j
\sim
\Normal\left(
0,
\frac{\sigma_0^2+1/N}{\omega_j}
\right),
\]
whereas
\[
x_{-j}\mid x_j,\omega_j,\gamma=j
\sim
\Normal_{J-1}\left(
x_j\mu_j,
\frac{c_j}{\omega_j}\bar S_j
\right)
\]
is unchanged. Let
\[
V_\beta=\sigma_0^2+\frac{1}{N},
\qquad
a=\frac{\nu_0+3}{2},
\qquad
b=\frac{\nu_0}{2}.
\]
Then
\begin{align}
p(x\mid \gamma=j)
=\;&
(2\pi)^{-J/2}|\bar S_j|^{-1/2}
V_\beta^{-1/2}
c_j^{-(J-1)/2}
\frac{b^a}{\Gamma(a)}
\notag\\
&\times
\int_0^\infty
\omega^{a+J/2-1}
\exp\left[
-\omega\left\{
b+\frac{x_j^2}{2V_\beta}
+\frac{q_R-x_j^2}{2c_j}
\right\}
\right]\;d\omega
\notag\\
=\;&
(2\pi)^{-J/2}|\bar S_j|^{-1/2}
V_\beta^{-1/2}
c_j^{-(J-1)/2}
\frac{b^a\Gamma(a+J/2)}{\Gamma(a)}
\notag\\
&\times
\left\{
b+\frac{x_j^2}{2V_\beta}
+\frac{q_R-x_j^2}{2c_j}
\right\}^{-(a+J/2)} .
\label{eq:common-noise-sqrt-scaled-marginal}
\end{align}
Since $|\bar S_j|=|\bar R|$ is common across SNPs, the PIP weights are
\begin{align}\label{eq:PIP-common-noise-mean-scaled}
w_j^{\mathrm{sqrt\text{-}scaled}}
=\;&
c_j^{-(J-1)/2}
\left\{
\frac{\nu_0}{2}
+\frac{x_j^2}{2(\sigma_0^2+1/N)}
+\frac{q_R-x_j^2}{2c_j}
\right\}^{-(\nu_0+J+3)/2},
\notag\\
\pi_j=\;&
\frac{w_j^{\mathrm{sqrt\text{-}scaled}}}
{\sum_{k=1}^Jw_k^{\mathrm{sqrt\text{-}scaled}}}.
\end{align}
This is better computationally: the integral over $\omega_j$ is now conjugate
and has a closed form. It is also closer to model A in the sense that both the
effective signal and the common noise share the same $1/\omega_j$ scale after
marginalizing over $\beta$. However, it still does not recover the exact
null-SNP irrelevance property, because the common statistic $q_R$ appears with
the SNP-specific coefficient $1/c_j$ inside the final bracket.


\section{Extension to susie}
Now we try to extend model B into a Susie-like model, where we allow more than
one non-zero effect. A natural extension is to keep the usual SuSiE additive
decomposition, but replace each single-effect update by the model B
single-effect calculation above.

\subsection{A first proposed multi-effect model}
Let there be $L$ single-effect components. For each component $\ell$, introduce
a causal index $\gamma_\ell$, an effect size $\beta_\ell$, and a relaxed LD
column $z_{\ell j}$ when $\gamma_\ell=j$. A direct extension of model B is
\[
x
\mid
\{\gamma_\ell,\beta_\ell,z_{\ell,\gamma_\ell},S_{\ell,\gamma_\ell}\}_{\ell=1}^L
\sim
\Normal_J\left(
\sum_{\ell=1}^L \beta_\ell u_{\gamma_\ell}(z_{\ell,\gamma_\ell}),
\frac{\bar R}{N}
\right),
\]
where $u_j(z_j)$ is the vector with entry $1$ at coordinate $j$ and entries
$z_j$ on the remaining coordinates. The priors are
\[
\gamma_\ell\sim\Unif\{1,\ldots,J\},
\qquad
\beta_\ell\sim\Normal(0,\sigma_0^2),
\]
and, conditional on $\gamma_\ell=j$,
\[
z_{\ell j}\mid S_{\ell j}
\sim
\Normal_{J-1}\left(\mu_j,\frac{S_{\ell j}}{\nu_0}\right),
\qquad
S_{\ell j}\sim\mathcal{IW}_{J-1}(\nu_0\bar S_j,\nu_0+J+1).
\]
This model gives each single-effect component its own relaxed active LD
column. This is the closest analogue of SuSiE: multiple effects add in the
mean, while each effect has a one-hot index and a scalar coefficient.

\paragraph{Sanity check: the $L=1$ case.}
This first proposed model is actually not the right extension of model B. The
issue is already visible when $L=1$. Take $\gamma=1$. Under the proposed model,
\[
x\mid \beta,z_1,S_1,\gamma=1
\sim
\Normal_J\left(
\beta\begin{pmatrix}1\\z_1\end{pmatrix},
\frac{\bar R}{N}
\right),
\]
so the covariance is fixed at $\bar R/N$ rather than using the
$(z_1,S_1)$-dependent covariance in model B. Using the block form
\[
\bar R=
\begin{pmatrix}
1 & \mu_1^\top\\
\mu_1 & \bar S_1+\mu_1\mu_1^\top
\end{pmatrix},
\]
we get
\[
x_1\mid \beta,\gamma=1
\sim
\Normal\left(\beta,\frac{1}{N}\right),
\]
and
\[
x_{-1}\mid x_1,\beta,z_1,\gamma=1
\sim
\Normal_{J-1}\left(
x_1\mu_1+\beta(z_1-\mu_1),
\frac{\bar S_1}{N}
\right).
\]
This is already different from model B, where the conditional mean is
$x_1z_1$ and does not contain $\beta$.

Now integrate over $z_1\mid S_1$. Since
\[
z_1\mid S_1\sim \Normal_{J-1}\left(\mu_1,\frac{S_1}{\nu_0}\right),
\]
we obtain
\[
x_{-1}\mid x_1,\beta,S_1,\gamma=1
\sim
\Normal_{J-1}\left(
x_1\mu_1,
\frac{\bar S_1}{N}+\frac{\beta^2}{\nu_0}S_1
\right).
\]
Therefore
\[
\begin{aligned}
p(x\mid \gamma=1)
=\;&
\Normal\left(x_1\mid 0,\frac{1}{N}+\sigma_0^2\right) \\
&\times
\E_{\beta\mid x_1}
\E_{S_1}
\left[
\Normal_{J-1}\left(
x_{-1}\mid x_1\mu_1,\,
\frac{\bar S_1}{N}+\frac{\beta^2}{\nu_0}S_1
\right)
\right],
\end{aligned}
\]
where
\[
\beta\mid x_1
\sim
\Normal\left(
\frac{N}{N+\sigma_0^{-2}}x_1,\,
\frac{1}{N+\sigma_0^{-2}}
\right),
\qquad
S_1\sim\mathcal{IW}_{J-1}(\nu_0\bar S_1,\nu_0+J+1).
\]
Thus the PIP weight under this proposed model is
\[
\boxed{
p(\gamma=j\mid x)
\propto
\Normal\left(x_j\mid 0,\frac{1}{N}+\sigma_0^2\right)
\E_{\beta\mid x_j}
\E_{S_j}
\left[
\Normal_{J-1}\left(
x_{-j}\mid x_j\mu_j,\,
\frac{\bar S_j}{N}+\frac{\beta^2}{\nu_0}S_j
\right)
\right].
}
\]
This is not the model B PIP. In particular, it depends on $x_{-j}$ through the
conditional residual $x_{-j}-x_j\mu_j$, and the integral is not conjugate
because the covariance is a sum of the fixed matrix $\bar S_j/N$ and the
random matrix $(\beta^2/\nu_0)S_j$. Therefore this first proposed
multi-effect generative model does not preserve the single-effect
irrelevance property we want.

\subsection{Variational approximation}
The failed $L=1$ check above means that we should not simply use the fixed
noise covariance $\bar R/N$ and then plug in the model B single-effect weights.
Still, it is useful to record what a SuSiE-like variational strategy would have
to look like. As in SuSiE, the exact posterior is not tractable when $L>1$
because the effects interact through the shared likelihood. A natural
mean-field approximation is
\[
q\left(
\{\gamma_\ell,\beta_\ell,z_{\ell,\gamma_\ell},S_{\ell,\gamma_\ell}\}_{\ell=1}^L
\right)
=
\prod_{\ell=1}^L
q_\ell(\gamma_\ell,\beta_\ell,z_{\ell,\gamma_\ell},S_{\ell,\gamma_\ell}).
\]
Each factor $q_\ell$ can be updated by holding the other components fixed in
expectation. Define the expected contribution of component $k$ by
\[
\hat b_k
=
\E_{q_k}\{\beta_k u_{\gamma_k}(z_{k,\gamma_k})\},
\]
and define the residual for updating component $\ell$ as
\[
x^{(\ell)}
=
x-\sum_{k\ne \ell}\hat b_k.
\]
The tempting coordinate update would be to apply the single-effect model B
formulas to $x^{(\ell)}$ in place of $x$. The sanity check above shows that
this update is not justified by the first proposed generative model with fixed
noise covariance $\bar R/N$. Instead, it should be viewed as a possible
pseudo-likelihood or variational working update that deliberately preserves the
model B single-effect behavior.

In particular, for each candidate SNP $j$, let $x_j^{(\ell)}$ be the $j$th
entry of the residual and let
\[
q_\ell^{R}
=
{x^{(\ell)}}^\top\bar R^{-1}x^{(\ell)}.
\]
The single-effect weight for component $\ell$ is
\[
w_{\ell j}
\propto
\left(
1+\frac{N{x_j^{(\ell)}}^2}{\nu_0}
\right)^{(\nu_0+3)/2}
\Normal\left(
x_j^{(\ell)}
\mid
0,\frac{1}{N}+\sigma_0^2
\right),
\]
and the variational PIP for component $\ell$ is
\[
\alpha_{\ell j}
=
\frac{w_{\ell j}}{\sum_{k=1}^J w_{\ell k}}.
\]
Conditional on $\gamma_\ell=j$, the variational posteriors for
$\beta_\ell,z_{\ell j},S_{\ell j}$ are exactly the model B posterior formulas
with $x$ replaced by $x^{(\ell)}$.

\subsection{Expected fitted effect}
For the SuSiE-style coordinate ascent algorithm, we also need the expected
fitted vector $\hat b_\ell$. Conditional on $\gamma_\ell=j$,
\[
\E(\beta_\ell\mid x^{(\ell)},\gamma_\ell=j)
=
\frac{N}{N+\sigma_0^{-2}}x_j^{(\ell)}.
\]
Also,
\[
\E(z_{\ell j}\mid x^{(\ell)},\gamma_\ell=j)
=
\frac{
\nu_0\mu_j+N x_j^{(\ell)}x_{-j}^{(\ell)}
}{
\nu_0+N{x_j^{(\ell)}}^2
}.
\]
Because $\beta_\ell$ and $(z_{\ell j},S_{\ell j})$ factorize conditional on
$x^{(\ell)}$ and $\gamma_\ell=j$ in the single-effect update, we can use
\[
\E\{\beta_\ell u_j(z_{\ell j})\mid x^{(\ell)},\gamma_\ell=j\}
=
m_{\beta,\ell j}
\begin{pmatrix}
1\\
m_{z,\ell j}
\end{pmatrix}_{j},
\]
where
\[
m_{\beta,\ell j}
=
\frac{N}{N+\sigma_0^{-2}}x_j^{(\ell)},
\qquad
m_{z,\ell j}
=
\frac{
\nu_0\mu_j+N x_j^{(\ell)}x_{-j}^{(\ell)}
}{
\nu_0+N{x_j^{(\ell)}}^2
}.
\]
Here the subscript $j$ on the vector means that the scalar $1$ is placed in
coordinate $j$ and $m_{z,\ell j}$ fills the remaining coordinates. Therefore
\[
\hat b_\ell
=
\sum_{j=1}^J
\alpha_{\ell j}
m_{\beta,\ell j}
\begin{pmatrix}
1\\
m_{z,\ell j}
\end{pmatrix}_{j}.
\]

\subsection{What is exact and what is approximate}
For $L=1$, this reduces exactly to model B. For $L>1$, the coordinate update is
a SuSiE-style variational approximation: each update treats the current
residual as the data for one single-effect model. This is not the exact
posterior under the full multi-effect model, but it preserves the key
single-effect behavior. In particular, the candidate weight $w_{\ell j}$
depends on the residual only through the marginal residual statistic
$x_j^{(\ell)}$, so the update retains the desired null-SNP irrelevance
property at the level of each single-effect regression.

\subsection{Hyperparameter learning}
The empirical Bayes updates for $\sigma_0^2$ and $\nu_0$ can also be lifted
from the single-effect case by averaging over components. For example, the
EM-style update for $\sigma_0^2$ becomes
\[
\sigma_{0,\mathrm{new}}^2
=
\frac{1}{L}
\sum_{\ell=1}^L\sum_{j=1}^J
\alpha_{\ell j}
\left(m_{\beta,\ell j}^2+v_\beta\right),
\qquad
v_\beta=\frac{1}{N+\sigma_{0,\mathrm{old}}^{-2}}.
\]
For $\nu_0$, one can use the same scalar EM objective as in the single-effect
case, replacing the sum over $j$ by a weighted sum over $(\ell,j)$ and using
the residual $x^{(\ell)}$ in the definitions of $q_R$, $x_j$, and the posterior
quantities. This gives a practical coordinate-ascent scheme: update the
single-effect factors one at a time, update fitted effects and residuals, and
periodically update $(\nu_0,\sigma_0^2)$ by empirical Bayes.

\section{Other models}
\subsection{Model C}
\begin{align}\label{eq:model-C}
\begin{cases}
x_1\mid \beta,\gamma=1
\sim
\Normal\left(\beta,\dfrac{1}{N}\right),\\
x_{-1}\mid x_1,z_1,\gamma=1
\sim
\Normal_{J-1}\left(
x_1 z_1,
\dfrac{1}{N} \bar S_1
\right),\\
z_1 | S_1 \sim \mathcal{N}\left(\mu_1, \dfrac{S_1}{\nu_0}\right),\\
S_1\sim\mathcal{IW}(\nu_0\bar S_1,\nu_0+J+1),\\
\beta\sim\Normal(0,\sigma_0^2).
\end{cases}
\end{align}
For a general causal SNP $j$, first integrate out $S_j$ from the prior on
$z_j$. Since
\[
z_j\mid S_j\sim \Normal_{J-1}\left(\mu_j,\frac{S_j}{\nu_0}\right),
\qquad
S_j\sim\mathcal{IW}(\nu_0\bar S_j,\nu_0+J+1),
\]
we have
\[
z_j\sim
t_{J-1,\nu_0+3}\left(
\mu_j,
\frac{1}{\nu_0+3}\bar S_j
\right).
\]
Equivalently, with the shape--rate convention,
\[
\lambda_j\sim\Gam\left(\frac{\nu_0+3}{2},\frac{1}{2}\right),
\qquad
z_j\mid \lambda_j
\sim
\Normal_{J-1}\left(
\mu_j,
\frac{1}{\lambda_j}\bar S_j
\right).
\]
Therefore, conditional on $x_j$ and $\lambda_j$,
\[
x_{-j}\mid x_j,\lambda_j,\gamma=j
\sim
\Normal_{J-1}\left(
x_j\mu_j,
\left\{
\frac{1}{N}
+\frac{x_j^2}{\lambda_j}
\right\}\bar S_j
\right).
\]
Let
\[
V_\beta=\sigma_0^2+\frac{1}{N},
\qquad
q_R=x^\top\bar R^{-1}x,
\qquad
r_j^2=q_R-x_j^2,
\]
where the last equality uses the block inverse identity
\[
r_j^2
=(x_{-j}-x_j\mu_j)^\top\bar S_j^{-1}(x_{-j}-x_j\mu_j).
\]
Then
\begin{align}
p(x\mid \gamma=j)
=\;&
\Normal(x_j;0,V_\beta)
(2\pi)^{-(J-1)/2}|\bar S_j|^{-1/2}
\notag\\
&\times
\int_0^\infty
\left\{
\frac{1}{N}
+\frac{x_j^2}{\lambda}
\right\}^{-(J-1)/2}
\notag\\
&\times
\exp\left[
-\frac{r_j^2}
{2\left\{1/N+x_j^2/\lambda\right\}}
\right]\notag\\
&\times
\frac{\left(1/2\right)^{(\nu_0+3)/2}}
{\Gamma((\nu_0+3)/2)}
\lambda^{(\nu_0+3)/2-1}
\exp\left\{-\frac{\lambda}{2}\right\}
\;d\lambda .
\label{eq:model-C-marginal}
\end{align}
Since $|\bar S_j|=|\bar R|$ is common across SNPs, the PIP weights are
\begin{align}\label{eq:PIP-model-C}
w_j^{C}
=\;&
\Normal(x_j;0,V_\beta)
\int_0^\infty
\left\{
\frac{1}{N}
+\frac{x_j^2}{\lambda}
\right\}^{-(J-1)/2}
\notag\\
&\times
\exp\left[
-\frac{q_R-x_j^2}
{2\left\{1/N+x_j^2/\lambda\right\}}
\right]
\lambda^{(\nu_0+3)/2-1}
\exp\left\{-\frac{\lambda}{2}\right\}
\;d\lambda,
\notag\\
\pi_j^C
=\;&
\frac{w_j^C}{\sum_{k=1}^J w_k^C}.
\end{align}
Thus model C has a one-dimensional integral per SNP. It does not preserve the
exact null-SNP irrelevance property: the common statistic
$q_R=x^\top\bar R^{-1}x$ appears inside an $x_j$-dependent integral, so it
cannot be factored out and cancelled in the PIP normalization.

\end{document}
